\documentclass[a4paper,12pt]{article}

% Paquetes básicos
\usepackage[utf8]{inputenc}
\usepackage[spanish]{babel}
\usepackage[T1]{fontenc}
\usepackage{geometry}
\usepackage{titlesec}
\usepackage{setspace}
\usepackage{hyperref}

% Márgenes y espaciado
\geometry{left=2.5cm,right=2.5cm,top=3cm,bottom=3cm}
\setstretch{1.3}

% Configuración de títulos
\titleformat{\section}{\normalfont\Large\bfseries}{\thesection}{1em}{}
\titleformat{\subsection}{\normalfont\large\bfseries}{\thesubsection}{1em}{}

% Documento
\begin{document}

% Portada
\begin{titlepage}
    \centering
    \vspace*{5cm}

    {\Huge \textbf{Informe del Trabajo Práctico Nº 2}\\[1.5cm]}

    {\Large Paradigmas de Programación - FIUBA \\[1.5cm]}

    \textbf{Integrantes:} \\[0.5cm]
    - Agustin Ferronato\\
    - Faustino Millares \\
    - Matias Llorens \\
    - Santiago Zottig \\[2cm]

    \vfill
    \textbf{Fecha:} \today

\end{titlepage}

% Índice (opcional)
\tableofcontents
\newpage

% Contenido
\section{Supuestos}
Aquí se detallan todos los supuestos tomados a partir del enunciado.
A continuación se detallan los supuestos que tomamos:

\begin{enumerate}
    \item Al repartir aleatoriamente las 10 cartas a cada jugador
     de su mazo, se le ofrece la posibilidad de descartar 2
     cartas por otras 2 del mazo, estas 2 cartas seran 
     descartadas aleatoriamente.
    \item Una carta con el modificador Agil al ser utilizada
    se ubicara aleatoriamente en una de las secciones de las 
    que se puede ubicar
    \item Una carta con el modificador de Medico al ser utilizada
    jugara automaticamente y de forma aleatoria una carta de la
    pila de descarte en caso de que esta contenga alguna
\end{enumerate}

\section{Diagrama de Clases}
Aquí se insertan o describen los diagramas de clases.

\section{Diagrama de Secuencia}
Aquí se inserta o explica el diagrama de secuencia.

\section{Diagrama de Paquetes}
Aquí se incluye el diagrama de paquetes.

\section{Detalles de Implementación}
Se describen detalles relevantes sobre la implementación, decisiones técnicas, patrones, etc.

\section{Excepciones}
Se enumeran las excepciones o casos especiales tratados en la implementación.

\end{document}
